%! Author = chtompki
%! Date = 1/14/26
\documentclass[12pt]{amsart}
% Default margins are too wide all the way around. I reset them here
\setlength{\hoffset}{-.75in}
\setlength{\textwidth}{6.5in}
\setlength{\voffset}{-.5in}
\setlength{\textheight}{9.0in}
\usepackage{amsmath}
\usepackage{leftindex}
\usepackage{amssymb}
\usepackage{amsthm}
\usepackage{enumitem}
\usepackage{setspace}

\newenvironment{solution}
{\begin{proof}[Solution]}
{\end{proof}}

\title{Combinatorics\\Exam 2}
\author{Rob Tompkins}
\date{\today}

\begin{document}
    \maketitle
    \date{\today}
    \begin{onehalfspace}
        % Problem 1
        \noindent\textbf{1.} (25pts): Let $a_n$ be the number of ways to partition an $n$-element
        set into ordered pairs (i.e. matchings where order within each pair matters).
        Find the exponential generating function for $a_n$.
        Derive a closed-form expression for $a_n$.

        \begin{solution}
            The pairs themselves are unordered, but the two entries in each
            pair are ordered.  On a fixed two-element set there are $2!$
            ordered pairs, so the exponential generating function for one
            pair is $2!x^2/2!=x^2$.  The exponential formula therefore gives
            \[
                A(x)=\sum_{n\ge0}a_n\frac{x^n}{n!}
                =\exp(x^2)=\sum_{m\ge0}\frac{x^{2m}}{m!}.
            \]
            Comparing coefficients yields
            \[
                a_{2m}=\frac{(2m)!}{m!},\qquad a_{2m+1}=0.
            \]
            Here $a_0=1$, corresponding to the empty partition.  Alternatively,
            arrange the $2m$ elements in a list and split it into consecutive
            pairs.  Each partition is counted $m!$ times, once for each
            ordering of its pairs, giving the same formula.
        \end{solution}

        % Problem 2
        \noindent\textbf{2.} (25pts): Let $a_n$ be the number of ways to distribute $n$ labled
        objects into any number of labeled boxes such that each box contains at least two
        objects.
        Construct the exponential generating function for $a_n$.
        Extract a formula for $a_n$.

        \begin{solution}
            Interpret ``any number of labeled boxes'' as summing over
            $k\ge0$, with the boxes labeled $1,\ldots,k$ for each fixed $k$.
            (Choosing arbitrary labels from an infinite supply would not give
            a finite count.)  A single box contains an unordered set of at
            least two labeled objects, with exponential generating function
            \[
                B(x)=\sum_{j\ge2}\frac{x^j}{j!}=e^x-1-x.
            \]
            For $k$ labeled boxes the generating function is $B(x)^k$.
            Consequently, as a formal power series,
            \[
                A(x)=\sum_{k\ge0}(e^x-1-x)^k
                    =\frac{1}{2+x-e^x}.
            \]
            If box $i$ contains $n_i$ objects, the number of distributions is
            the multinomial coefficient $n!/(n_1!\cdots n_k!)$.  Thus an
            explicit finite formula is
            \[
                a_n=\sum_{k=0}^{\lfloor n/2\rfloor}
                \ \sum_{\substack{n_1+\cdots+n_k=n\\n_i\ge2}}
                \frac{n!}{n_1!\cdots n_k!}.
            \]
            The $k=0$ term is $1$ when $n=0$ and $0$ otherwise.
            Equivalently, $a_n=\sum_k k!\left\{\begin{smallmatrix}n\\k
            \end{smallmatrix}\right\}_{\ge2}$, where the symbol counts
            partitions into $k$ blocks, each of size at least two.
        \end{solution}

        % Problem 3
        \noindent\textbf{3.} (25pts): The Catalan nunbers are defined by $C_n =
        \frac{1}{1+n}\binom{2n}{n}$.
        Probe that $C_n$ satisfies the recurrence $C_n=\sum_{k=0}^{n-1}C_k C_{n-1-k}$.
        Give a combinatorial interpretation of this recurrence.

        \begin{proof}
            Let $C(x)=\sum_{n\ge0}C_nx^n$.  The formal binomial expansion gives
            \[
                \frac{1-\sqrt{1-4x}}{2x}
                =\sum_{n\ge0}\left[-\frac12
                                 \binom{1/2}{n+1}(-4)^{n+1}\right]x^n
                =\sum_{n\ge0}\frac{1}{n+1}\binom{2n}{n}x^n.
            \]
            Hence $C(x)=(1-\sqrt{1-4x})/(2x)$, which satisfies
            $C(x)=1+xC(x)^2$.  Taking the coefficient of $x^n$, for $n\ge1$,
            proves
            \[
                C_n=\sum_{k=0}^{n-1}C_kC_{n-1-k},\qquad C_0=1.
            \]
            Combinatorially, $C_n$ counts Dyck paths: paths with $n$ up-steps
            and $n$ down-steps that never go below height zero.  Indeed, the
            reflection principle counts the paths that go below zero as
            $\binom{2n}{n+1}$, so the number of Dyck paths is
            $\binom{2n}{n}-\binom{2n}{n+1}=C_n$.
            Every nonempty Dyck path decomposes uniquely as $U P D Q$, where
            the displayed $D$ is the first return to height zero.  If $P$ has
            semilength $k$, then $Q$ has semilength $n-1-k$.  There are
            $C_kC_{n-1-k}$ choices for these two paths; summing over $k$
            gives the recurrence.
        \end{proof}

        % Problem 4
        \noindent\textbf{4.} (25 pts): How many integers between 1 and 1000 (inclusive) are
        divisible by at least one of 3,5, or 7?
        Use the Inclusion-Exclusion Principle and show all steps.

        \begin{solution}
            Let $A_d=\{m\in\{1,\ldots,1000\}:d\mid m\}$.
            Since $3,5,7$ are pairwise coprime, their intersections consist
            of multiples of the corresponding products.  Thus
            \begin{align*}
                |A_3|&=333,& |A_5|&=200,& |A_7|&=142,\\
                |A_3\cap A_5|&=\lfloor1000/15\rfloor=66,\\
                |A_3\cap A_7|&=\lfloor1000/21\rfloor=47,\\
                |A_5\cap A_7|&=\lfloor1000/35\rfloor=28,\\
                |A_3\cap A_5\cap A_7|&=\lfloor1000/105\rfloor=9.
            \end{align*}
            By inclusion--exclusion,
            \begin{align*}
                |A_3\cup A_5\cup A_7|
                &=333+200+142-66-47-28+9\\
                &=\boxed{543}.
            \end{align*}
        \end{solution}

        % Problem 5
        \noindent\textbf{5.} (25pts): Suppose we have an infinite supply of pennies and we would
        like to stack them up in rows.
        Find the number of different arrangements possible subject to the following conditions:
        \begin{itemize}
            \item[(i)] the bottom row consists of $n$ pennies, each touching its two neighbors,
            except the ends where $n\ge 1$;
            \item[(ii)] a soin that does not belong to a row sits on two coins below it;
            \item[(iii)] do not distinguish between heads and tails.
        \end{itemize}

        \begin{solution}
            Interpret condition (ii) as saying that every coin outside the
            \emph{bottom} row rests on two adjacent coins immediately below
            it.  Upper rows may have gaps, and left--right reflections are
            counted separately.  Let $f_n$ count these stacks on a fixed
            bottom row of $n$ coins, and set $f_0=1$.

            There are $n-1$ possible coin positions immediately above the
            bottom row.  If this entire next row is occupied, the choices
            above it form a stack with bottom row of length $n-1$, giving
            $f_{n-1}$ arrangements.

            Otherwise, let $k$ be the first vacant position in that row,
            counted from the left, where $1\le k\le n-1$.  The first $k-1$
            positions are occupied, and the choices above these coins give
            $f_{k-1}$ stacks.  To the right of the vacancy there is an arbitrary
            stack on the last $n-k$ bottom coins, giving $f_{n-k}$ choices.
            No coin can bridge the vacancy: any such coin would require,
            through its supports, the missing coin.  The two choices are
            therefore independent.  Consequently,
            \[
                f_n=f_{n-1}+\sum_{k=1}^{n-1}f_{k-1}f_{n-k}
                =\sum_{k=1}^{n}f_{k-1}f_{n-k}.
            \]
            This is the Catalan recurrence with $f_0=1$, so
            \[
                f_n=C_n=\frac{1}{n+1}\binom{2n}{n}.
            \]
            For example, $f_1=1$, $f_2=2$, and $f_3=5$.
        \end{solution}


        % Problem 6
        \noindent\textbf{6.} (25pts): Consider the colorings of the vertices of a regular hexagon using
        3 colors.
        Determine the cycle index polynomial of the dihedral group $D_6$.
        Use Polya's Enumeration Theorem to compute the number of distinct colorings.

        \begin{solution}
            Write $s_j$ for a cycle of length $j$.  The group $D_6$ has
            $12$ elements, with the following cycle types on the six vertices:
            \[
                \begin{array}{c|c|c}
                    \text{Symmetry}&\text{Number}&\text{Cycle monomial}\\ \hline
                    \text{Identity}&1&s_1^6\\
                    \text{Rotations by }\pm60^\circ&2&s_6\\
                    \text{Rotations by }\pm120^\circ&2&s_3^2\\
                    \text{Rotation by }180^\circ&1&s_2^3\\
                    \text{Reflections through opposite vertices}&3&s_1^2s_2^2\\
                    \text{Reflections through opposite edges}&3&s_2^3
                \end{array}
            \]
            Averaging these monomials gives
            \[
                Z(D_6)=\frac1{12}
                    (s_1^6+2s_6+2s_3^2+4s_2^3+3s_1^2s_2^2).
            \]
            Each cycle must be monochromatic in a coloring fixed by its
            permutation.  With three available colors, Polya's theorem
            therefore gives
            \[
                Z(D_6)(3,3,3,3,3,3)
                =\frac{3^6+2(3)+2(3^2)+4(3^3)+3(3^4)}{12}
                =\frac{1104}{12}=\boxed{92}.
            \]
            Colors are distinct, and there is no requirement to use all three.
        \end{solution}

        % Problem 7
        \noindent\textbf{7.} (25pts): Modle a benzene ring as a hexagon.
        Suppose we substitute hydrogen atoms with two types of atoms ($A$ and $B$).
        Using Polya's Enumeration Theorem, determine the number of distinct substitution patterns
        (up to rotational and reflectional symmetry).

        \begin{solution}
            If every site is occupied by either $A$ or $B$, these are
            two-colorings of a hexagon under $D_6$.  Substituting $s_j=2$
            into the cycle index from Problem 6 gives
            \[
                Z(D_6)(2,2,2,2,2,2)
                    =\frac{2^6+2(2)+2(2^2)+4(2^3)+3(2^4)}{12}
                    =13.
            \]
            More precisely, substituting $s_j=a^j+b^j$ gives the pattern
            inventory
            \begin{align*}
                P(a,b)=\frac1{12}\bigl(& (a+b)^6+2(a^6+b^6)
                +2(a^3+b^3)^2\\
                &+4(a^2+b^2)^3+3(a+b)^2(a^2+b^2)^2\bigr)\\
                ={}&a^6+a^5b+3a^4b^2+3a^3b^3
                +3a^2b^4+ab^5+b^6.
            \end{align*}
            Thus, for example, exactly three $A$ atoms and three $B$ atoms
            give three patterns.  If both types must occur, the total is
            $13-2=11$.

            If the wording instead allows some hydrogens to remain, each
            site has three states, $H,A,B$.  The answer is then $92$,
            as in Problem 6, including the unsubstituted ring; requiring
            at least one substitution gives $92-1=91$ patterns.
        \end{solution}

        % Problem 8
        \noindent\textbf{8.} (25pts): Show that: $R(3, 5) > 13$ and $R(3,5)\le 14$.

        \begin{proof}
            Recall that $R(3,5)$ is the least $N$ such that every graph on
            $N$ vertices contains a triangle or an independent set of size $5$.

            \noindent\emph{Lower bound.}
            Let $G$ have vertex set $\mathbb Z/13\mathbb Z$, with $u$ adjacent
            to $v$ exactly when $u-v\in\{1,-1,5,-5\}$ modulo $13$.
            This graph is triangle-free.  Indeed, the neighbors of $0$ are
            $1,5,8,12$, and the positive differences between distinct members
            are $3,4,7,11$; none is congruent to $\pm1$ or $\pm5$.
            Thus no two neighbors of $0$ are adjacent, and translation
            proves the same statement at every vertex.

            Suppose five vertices were independent.  List them in cyclic
            order, and let $g_1,\ldots,g_5$ be their positive cyclic gaps.
            No gap is $1$, and their sum is $13$.  Thus, up to reordering,
            the only possibilities are
            \[
                (5,2,2,2,2),\qquad(4,3,2,2,2),\qquad(3,3,3,2,2).
            \]
            The first possibility contains a gap of $5$, which is an edge.
            In the second, the gap $3$ must have a neighboring gap $2$,
            since there is only one gap $4$.  In the third, a gap $2$ and a
            gap $3$ must be consecutive somewhere around the cycle.
            In either case, two consecutive gaps sum to $5$, again giving
            an edge between two selected vertices.  This contradiction
            proves that $G$ has no independent set of size $5$.  Coloring
            its edges red and its nonedges blue gives no red triangle and
            no blue $K_5$, so $R(3,5)>13$.

            \noindent\emph{Upper bound.}
            First, every graph on six vertices has a triangle or an
            independent triple.  To see this, color edges red and nonedges
            blue.  A vertex has at least three incident edges of one color.
            Among their other endpoints, an edge of that color completes a
            monochromatic triangle; if none exists, the endpoints themselves
            form a triangle of the other color.  Thus $R(3,3)\le6$.

            Next, a triangle-free graph on nine vertices must have an
            independent set of size $4$.  Otherwise every vertex has degree
            at most $3$, since its neighbors are independent.  If a vertex
            had degree at most $2$, it would have at least six other
            nonneighbors; by the preceding result these contain an
            independent triple, which together with the vertex gives an
            independent set of size $4$.  All degrees would therefore be
            $3$, impossible because their sum would be $27$, an odd number.
            This proves $R(3,4)\le9$.

            Finally, let $G$ be triangle-free on $14$ vertices, and choose a
            vertex $v$.  If $\deg(v)\ge5$, five of its neighbors are an
            independent set.  Otherwise $v$ has at least $13-4=9$ other
            nonneighbors.  Their induced graph is triangle-free, so it
            contains an independent set of size $4$; adjoining $v$ gives
            an independent set of size $5$.  Hence $R(3,5)\le14$, and
            together the two bounds yield
            \[
                R(3,5)=14.
            \]
        \end{proof}

    \end{onehalfspace}
\end{document}
